\chapter{Aspect}
\sloppy

\begin{sourcebox}
Although the difference in meaning between the imperfective and perfective aspects is of little importance in mathematics, their difference in form provides the natural framework for a systematic study of Russian vocabulary.
\end{sourcebox}

\section{Difference in meaning between the two aspects}

As we have seen in Chapter III, a given Russian verb is imperfective (e.g. \ru{чита́ть} \emph{to read}) or perfective (e.g. \ru{прочита́ть} \emph{to read}) according as
\begin{enumerate}[label=\arabic*)]
\item (imperfective): it refers to the occurrence of a state or action in the present, past or future, with no implication about its completion;
\item (perfective): it implies something about completion in the past or future (so that, in particular, a perfective verb cannot be in the present tense).
\end{enumerate}

The difference in meaning between a past imperfective and a past perfective is sometimes indicated in English by the word-order; e.g., \emph{we have already done that job} (i.e. many times before; past imperfective in Russian) and \emph{we already have that job done} (i.e., now that we've got that job completed, let's go on to another; past perfective in Russian). For the two Russian verbs, e.g. \ru{чита́ть} and \ru{прочита́ть} \emph{to read}, of an aspect-pair we shall sometimes use the symbol $\sim$, e.g. \ru{чита́ть $\sim$ прочита́ть}, the imperfective being put first.

For the difference between future imperfective and future perfective, compare
\begin{center}
\begin{tabular}{@{}>{\cyrillicfont}p{68mm}p{70mm}@{}}
бу́дем получа́ть не́которые результа́ты & we shall obtain some results\\
\normalfont(\ru{получа́ть} \emph{to obtain}, imperfective) & \\
\addlinespace
полу́чим оконча́тельный результа́т & we shall obtain a definitive result\\
\normalfont(\ru{получи́ть} \emph{to obtain}, perfective) &
\end{tabular}
\end{center}

The two verbs of an aspect-pair do not differ in lexical meaning; i.e., \ru{чита́ть} and \ru{прочита́ть} both mean \emph{to read}, but their aspectual difference carries many subtle implications that can only be paraphrased in English. These implications permeate all Russian literature and conversation, sometimes with humorous or poetic effect, but they are of little interest in mathematics. In a sentence like \emph{if we take such-and-such steps, we obtain (or shall obtain) such-and-such a result}, the present imperfective \ru{получа́ем} and the future perfective \ru{полу́чим} are used interchangeably, the first one meaning \emph{while we are taking the steps, we are obtaining the result} and the second \emph{when we have taken the steps, we shall have completed the obtaining of the result}, but this distinction is of no more mathematical interest than the corresponding distinction in English between \emph{we obtain} and \emph{we shall obtain}.

So the question arises: Why then is it necessary for us to study aspect in detail?

To this question there are two answers. The first reason is that examples of both aspects occur so frequently in mathematics that the reader must know how they are constructed in order not to be confused by a maze of unfamiliar forms. But there is another reason, of such basic importance that it provides the framework for the whole book: namely, the entire learned vocabulary of Russian (and to a lesser extent, English) is closely related to the compounding of verbs with prepositions, which in turn is closely related to the construction of perfective verbs, in the following way. In spite of the fact that other methods of forming the perfective were typical for Indo-European and Proto-Slavic, and are still fairly important in modern Russian (see §6), the prefixing of prepositions (as e.g. in \ru{писа́ть $\sim$ написа́ть} \emph{to write}) is the regular method in modern Russian. A compound like \ru{написа́ть} \emph{to write}, where the prefix \ru{на-} merely indicates the aspect without making any other change in the meaning of the simple verb \ru{писа́ть} \emph{to write}, is called an \emph{aspectual compound}, whereas a compound like \ru{вписа́ть} \emph{to inscribe}, in which the prefix \ru{в-} \emph{in} produces a definite change of meaning, is called a \emph{lexical compound}, and the two types of compound verbs, together with their derived nouns and adjectives, account for almost all the Russian vocabulary of practical interest to a mathematician.

\section{Aspect regarded as a correspondence}

Aspect may be described, in set-theoretical terms, as a partitioning of the class of all Russian verbs into two (almost) mutually exclusive subclasses with an (almost) one-to-one correspondence between them; i.e. almost every verb has its uniquely determined aspect-partner, with the same lexical meaning.

However, there are exceptions of many kinds. The two classes are not quite mutually exclusive, since several imperfective verbs are also their own perfectives; e.g. foreign verbs in \ru{-ировать}: \ru{дифференци́ровать} \emph{to differentiate}, \ru{регистри́ровать} \emph{to register}, \ru{редакти́ровать} \emph{to edit}. Nor is the imperfective-perfective correspondence always one-to-one. It may be one-to-two (or more), since a given imperfective verb may form its perfective with two or more prepositional prefixes, corresponding to slightly different meanings of the simple verb; e.g. \ru{учи́ть} means either \emph{to learn}, when the perfectivizing prefix is \ru{вы-}, or \emph{to teach}, when it is \ru{вы-, на-} or \ru{об-} corresponding to slightly different kinds of teaching; and the foreign verbs just mentioned have as perfective partners not only themselves but (for emphasis on the perfectivization) \ru{продифференци́ровать}, \ru{зарегистри́ровать}, and \ru{отредакти́ровать}. Also, the correspondence may be one-to-zero; e.g., from their very meaning the imperfectives \ru{иска́ть} \emph{to seek} and \ru{зави́сеть} \emph{to depend on} have no perfective partners; and finally the correspondence may be zero-to-one; e.g. the perfective \ru{пона́добиться} \emph{to be necessary} (i.e. to have become necessary) has no imperfective partner.

\section{Perfective partners in contrast to lexical compounds}

The above examples show that the commonest method (for exceptions see §6.2) of forming the perfective partner for a given imperfective verb is to prefix a particular preposition, different from verb to verb. So it is natural to ask two questions:
\begin{enumerate}[label=\roman*)]
\item how is the preposition chosen for a given verb?
\item what is the logical connection between prepositions and the idea of completeness characteristic of perfective verbs?
\end{enumerate}

To answer these questions, let us examine the simple (i.e. uncompounded) verb \ru{писа́ть} \emph{to write} with the aspectual compound (i.e. perfective partner) \ru{написа́ть} and with many lexical compounds (i.e. compounds meaning something else than just \emph{to write}) e.g. \ru{описа́ть} \emph{to describe} (to write about). In view of the fact that the two verbs of an aspect-pair differ only in aspect, it is clear that \ru{описа́ть} \emph{to describe} cannot be a perfective partner of \ru{писа́ть} \emph{to write}, since for this particular verb the preposition \ru{о} \emph{about} produces too great a change of lexical meaning, from \emph{to write} to \emph{to describe}. But \ru{на} \emph{on} does not add very much to the lexical meaning of \ru{писа́ть} \emph{to write}, since all writing must obviously be on something, e.g. paper or stone, and must be completed by the time it is actually on the stone.

On the other hand, the prefix \ru{о-} can very well be used to form the perfective of a verb like \ru{характеризова́ть} \emph{to characterize}, since the uncompounded verb already contains the idea of \emph{about} (to characterize is to describe, i.e. to write about), which the prefix \ru{о-} only emphasizes, so as to suggest perfection or completion, although for many verbs the choice of perfectivizing prepositions will be less clear-cut.

In answer to the second question: namely, what is the connection between prepositions and the idea of completeness? let us note that most of the 19 prefixes listed in Chapter III as being available to form a compound verb, namely the 15 proper prepositions \ru{в, до, за}, etc. and the four inseparable prefixes \ru{воз-/вз-, вы-, пере-, раз-}, naturally suggest completion by their inherent meaning; e.g. \ru{за-, пере-, про-} suggest through-and-through; \ru{из-, от-, с-, у-} suggest out-and-out; and \ru{воз-/вз-, вы-, до-} give the same effect of completion as the \emph{up} in the English \emph{to finish up}. Consequently, any compound verb (e.g. \ru{описа́ть} \emph{to describe}) is perfective, unless the prefix is counteracted by one of the imperfectivizing suffixes discussed in the next section (for a few exceptions like \ru{получа́ть} \emph{obtain}, impf., see §6.2). However, although all these 19 prefixes form lexical compounds, only 15 of them can be used as mere perfectivizers (i.e. to form an aspectual compound), since the sharply local meaning of four of them \ru{в} \emph{into}, \ru{до} \emph{up to}, \ru{над} \emph{over} and \ru{пред} (or \ru{перед}) \emph{in front of} makes too great a change in the meaning of the simple verb; e.g. \ru{писа́ть} \emph{to write} becomes \ru{вписа́ть} \emph{to inscribe}. So there remain exactly 15 possible perfectivizers (the 4 inseparable prefixes and 11 of the 15 proper prepositions), some of which serve for many verbs and some for very few; a local preposition like \ru{на} with the meaning \emph{on}, can be the perfectivizer for only a few verbs, but the vague meaning of \ru{по} \emph{according to} makes it by far the commonest perfectivizer of the fifteen.\corrnote{原文在排除 19 个前缀中的 4 个之后仍写成 “there remain exactly 19 possible perfectivizers”，且括号内仍写 “15 of the 19 proper prepositions”。按原文自己的计数应为 15，即 4 个 inseparable prefixes 加 11 个可作 mere perfectivizers 的 proper prepositions；此处正文作最小更正。}

\section{The imperfectivizing suffixes}

If \ru{написа́ть} is the perfective of \ru{писа́ть}, then conversely the imperfective of \ru{написа́ть} is \ru{писа́ть}. But the simple verb \ru{писа́ть} \emph{to write} is also found in (lexical) compounds with many other prefixes, e.g. \ru{в-} \emph{into} (Latin \emph{in}), \ru{о-} \emph{about} (Latin \emph{circum} or \emph{de}), \ru{пере-} \emph{through} (Latin \emph{trans}), thereby forming (perfective) loan-translations like
\begin{center}
\small
\begin{tabular}{@{}>{\cyrillicfont}l p{105mm}@{}}
вписа́ть & \emph{inscribe}; e.g. a circle in a triangle\\
записа́ть & \emph{rescribe}; e.g. rewrite a formula in a different way (cf. \ru{за́пись} notation)\\
переписа́ть & \emph{transcribe}; e.g. make a list of, index\\
подписа́ть & \emph{subscribe}; e.g. write as a subscript\\
приписа́ть & \emph{ascribe}, assign to
\end{tabular}
\end{center}

How then do we form the imperfectives of these new perfective verbs?

As is shown by the table in the next section, there are two basic methods:
\begin{enumerate}[label=\arabic*)]
\item infinitives not ending in \ru{-ить} insert the imperfectivizing suffix \ru{-ыва-}, e.g. \ru{описа́ть} \emph{describe} forms the imperfective \ru{опи́сывать}. The \ru{-ыва-} is replaced by \ru{-ива-} where demanded by the spelling-rule, and in a few other verbs as well, e.g. \ru{рассмотре́ть} forms the imperfective \ru{рассма́тривать} (for these latter verbs the vowel gradation from \ru{о} to \ru{а} is typical);
\item infinitives in \ru{-ить} (except a few, e.g. \ru{перестро́ить} \emph{reconstruct} forms the imperfective \ru{перестра́ивать}) change the \ru{-ить} to \ru{-ять} (with inserted \ru{-л-} after \ru{в, м, п} or \ru{ф}; e.g. \ru{вставля́ть $\sim$ вста́вить}).
\end{enumerate}

Column 10 of the table in the next section provides the following examples: for the first method, \ru{писа́ть} (rows 1, 8, 9, 13, 14, 15), \ru{заде́рживать} (row 5), \ru{изы́скивать} (row 6), \ru{проде́лывать} (row 16), \ru{рассма́тривать} (row 17), \ru{скрыва́ть} (row 18); and for the second, \ru{возбужда́ть} (row 2; note \ru{-ать} for \ru{-ять} after \ru{ж} by the spelling-rule), \ru{выделя́ть} (row 3), \ru{доверя́ть} (row 4), \ru{направля́ть} (row 7), \ru{отделя́ть} (row 10), and \ru{уклоня́ться} (row 19).

\section{Table of aspectual and lexical compounds of verbs}

The entries in the following Table should be pronounced aloud, with the idea in mind that others like them will turn up in Chapter V. Proceed across each row, verifying that a given form actually belongs in the given position, with respect to both prefix and ending. In particular, the Table should be regarded as an opportunity to become familiar with the English (i.e. Latin) equivalents of the 19 prefixes.

Taking Row 9 as an example, the structure of the table is as follows. Column 1 lists the Russian prepositional prefix \ru{о-} \emph{about}, Column (2) lists the corresponding Latin prefix \emph{circum}, etc., (3) the English meaning of the prefix, (4) an example of an imperfective verb \ru{характеризова́ть} perfectivized with this prefix (5), followed by the English meaning (6) of the aspect-pair of columns (4) and (5). Then Column (7) gives an example of an imperfective verb \ru{писа́ть} (with its perfective \ru{написа́ть} in parentheses below it) which forms the lexical compound \ru{описа́ть} in Column (9), followed by its imperfective partner \ru{опи́сывать} in Column (10) and their meaning in Column (11). Finally, Column (12) gives the corresponding Russian roots; e.g. \ru{характеризова́ть} is a foreign borrowing, as indicated by the dash, and \ru{писа́ть} has the root \ru{пис-}.

Note that \ru{воз-, из-, раз-} become \ru{вос-, ис-, рас-} before an unvoiced consonant (see p. 26), and \ru{о-, от-, под-, с-} become \ru{обо-, ото-, подо-, со-} before certain double consonants.

\begingroup
\footnotesize
\setlength{\tabcolsep}{2.5pt}
\renewcommand{\arraystretch}{1.08}
\begin{longtable}{@{}r >{\cyrillicfont}p{19mm} p{17mm} p{22mm} >{\cyrillicfont}p{32mm} p{29mm}@{}}
\toprule
& \normalfont Russian prefix & Latin prefix & Meaning & \normalfont Example of an imperfective verb / perfective partner & Meaning of aspect-pair\\
\midrule
\endhead
1 & в- & in & in, into & --- & ---\\
2 & воз-/вз- & ex & out, up & пользоваться $\sim$ воспользоваться & use, exploit\\
3 & вы- & ex & out, up & учить $\sim$ выучить & teach\\
4 & до- & sub & up to & --- & ---\\
5 & за- & de & beyond & регистрировать $\sim$ зарегистрировать & register\\
6 & из- & ex & out, from & пробовать $\sim$ испробовать & test, probe\\
7 & на- & in, di- & on, over & писать $\sim$ написать & write\\
8 & над- & super & over & --- & ---\\
9 & о- & de, circum & about & характеризовать $\sim$ охарактеризовать & characterize\\
10 & от- & dis-, se-, re- & out & редактировать $\sim$ отредактировать & edit\\
11 & пере- (or пре-) & trans & across & путать $\sim$ перепутать & entangle\\
12 & по- & --- & according to & строить $\sim$ построить & construct\\
13 & под- & sub & under & готовить $\sim$ подготовить & prepare\\
14 & пред- & pre & before & --- & ---\\
15 & при- & ad & to, at & близиться $\sim$ приблизиться & approach\\
16 & про- & pro, per & about, through & диф\-ференци́ровать $\sim$ про\-диф\-ференци́ровать & differentiate\\
17 & раз- & dis & apart & делить $\sim$ разделить & divide\\
18 & с- & con- & with, from & делать $\sim$ сделать & do\\
19 & у- & dis-, de- & near, away & множить $\sim$ умножить & multiply\\
\bottomrule
\end{longtable}
\endgroup

\begingroup
\footnotesize
\setlength{\tabcolsep}{2.2pt}
\renewcommand{\arraystretch}{1.08}
\begin{longtable}{@{}r >{\cyrillicfont}p{27mm} p{16mm} >{\cyrillicfont}p{21mm} >{\cyrillicfont}p{24mm} p{23mm} >{\cyrillicfont}p{18mm}@{}}
\toprule
& \normalfont Example of an imperfective verb (perfective in parentheses) & Meaning & \normalfont Lexical compound & \normalfont Imperfective of lexical compound & Meaning & \normalfont Russian roots\\
\midrule
\endhead
1 & писать (написать) & write & вписать & вписывать & inscribe & ---, пис-\\
2 & будить (пробудить) & awaken & возбудить & возбуждать & excite & польз-, буд-\\
3 & делить (разделить) & divide & выделить & выделять & extract & уч-, дел-\\
4 & верить (поверить) & believe & доверить & доверять & entrust to & ---, вер-\\
5 & держать (no pf.) & hold & задержать & задерживать & detain & ---, держ-\\
6 & искать (no pf.) & seek & изыскать & изыскивать & examine & ---, иск-\\
7 & править (no pf.) & rectify & направить & направлять & direct & пис-, прав-\\
8 & писать (написать) & write & надписать & надписывать & superscribe & ---, пис-\\
9 & писать (написать) & write & описать & описывать & describe & ---, пис-\\
10 & делить (разделить) & divide & отделить & отделять & select & ---, дел-\\
11 & писать (написать) & write & переписать & переписывать & transcribe & пут-, пис-\\
12 & --- & --- & \multicolumn{4}{l}{\normalfont (almost all compounds with \ru{по-} are purely aspectual)}\\
13 & писать (написать) & write & подписать & подписывать & subscribe to & готов-, пис-\\
14 & писать (написать) & write & предписать & предписывать & prescribe & ---, пис-\\
15 & писать (написать) & write & приписать & приписывать & adscribe, assign to & близ-, пис-\\
16 & делать (сделать) & do & проделать & проделывать & perform & ---, дел-\\
17 & смотреть (посмотреть) & look at & рассмотреть & рассматривать & discuss, consider & дел-, смотр-\\
18 & крыть (покрыть) & cover & скрыть & скрывать & conceal & дел-, кры-\\
19 & клониться (no pf.) & bend, incline & уклониться & уклоняться & digress, deviate & мног-, клон-\\
\bottomrule
\end{longtable}
\endgroup

\section{Notes on the Table}

\begin{enumerate}[label=\arabic*)]
\item The Russian prefixes \ru{до, за, от, про, с, у} do not correspond very closely to any fixed Latin preposition, so that the loan-translations formed with them come from Latin words beginning with various prefixes; e.g. \ru{отделя́ть} \emph{select}, \ru{относи́ть} \emph{refer}, etc.

\item Some of the Russian verbs discussed in Chapter V fail to fit into the table because they construct their perfective partners otherwise than by prefixing a preposition; in the pair \ru{дава́ть, дать} \emph{to give} it is clear that the perfective is the older of the two forms (\ru{-ва-} is an imperfectivizing suffix found in a few other verbs as well); in \ru{бра́ть, взять} \emph{to take}, the two aspect-partners are constructed from different roots (\ru{бер-} = Lat. \emph{-fer} \emph{take}; and \ru{-я-} \emph{take} with the perfectivizing prefix \ru{вз-}), and similarly for \ru{говори́ть, сказа́ть} \emph{to say} (cf. verbs in English like \emph{go, went}). In \ru{получа́ть, получи́ть} \emph{to obtain} (with the root \ru{луч-} \emph{admit, receive}, from which comes also the important noun \ru{слу́чай} \emph{event, instance, case}) the unusual presence of a prefix in the imperfective results from confusion of forms borrowed from Old Church Slavonic.

\item Some lexical compounds are loan-translated from Latin verbs with two prefixes, e.g. \ru{подраздели́ть} \emph{to subdivide} (\ru{под-} sub-, \ru{раз-} di-); \ru{совпада́ть} \emph{to coincide} (\ru{со-} co-, \ru{в-} in-, \ru{пад-} fall, see p. 4). Others have two prefixes even though the Latin verb has only one; e.g. \ru{производи́ть} \emph{produce}.

\item No exercises are given here, since the discussion of aspect is intended only as a preparation for the vocabulary of the next chapter, where we shall find that most verbs, together with their derived nouns and adjectives, fit very well into the framework of the present chapter. But as might be expected, the very commonest verbs (see e.g. §§2, 3, 4 of Chapter V) are also the most irregular. For example, the verb \ru{быть} (no pf.) \emph{to be} (perhaps the most irregular of all verbs, in English as well as Russian) fails entirely to fit into the scheme, with forms coming from at least two different roots; on the one hand, words like \ru{быть} \emph{to be}, \ru{был} \emph{was}, and their derivatives \ru{быва́ть} \emph{to occur often}, \ru{бы́вший} \emph{the past} etc., and on the other hand, \ru{есть} \emph{is}, \ru{суть} \emph{are}, with derivatives like \ru{есте́ственный} \emph{natural}, \ru{существова́ть} \emph{exist}, \ru{существова́ние} \emph{existence}, \ru{суще́ственно} \emph{essentially}, \ru{су́щий} \emph{essential}, \ru{осуществля́ть} \emph{realize}, \ru{отсу́тствие} \emph{absence}, \ru{отсу́тствовать} \emph{be absent}, etc.
\end{enumerate}

\fussy
